% =====================================================================
\chapter{Plano de desenvolvimento, arquitetura e requisitos}\label{cap:arquitetura}
% =====================================================================

O problema foi organizado com o método 5W2H, já usado pelo grupo: o quê (receber e guardar leituras que geram alertas, com garantias de dado crítico), por quê (uma leitura perdida ou adulterada atrasa um alerta), onde (estações em campo, servidor da Defesa Civil e navegador dos operadores), quando (uma leitura por minuto e um lote a cada 10 minutos, sem parar), quem (dispositivos, operadores, órgãos de consulta e administrador), como (PostgreSQL com regras e segurança internas, servidor Node.js e painel web) e quanto (só software livre). Em seguida, a primeira versão do banco foi executada do zero e submetida a casos de borda; o \autoref{qua:diagnostico} resume os problemas encontrados, que orientaram as melhorias. Cada melhoria foi feita numa ramificação própria no GitHub e testada recriando o banco do zero.

\begin{quadro}[htb]
  \caption{Problemas encontrados na primeira versão}\label{qua:diagnostico}
  \footnotesize
  \begin{tabularx}{\textwidth}{|L{5.9cm}|Y|}
    \hline
    \textbf{Problema observado} & \textbf{Risco para dados críticos} \\ \hline
    Partições criadas só na carga, até dezembro de 2026 & Toda transmissão falharia a partir de janeiro de 2027 \\ \hline
    Uma leitura com data inválida derrubava o lote inteiro & Uma leitura ruim descartava as outras nove \\ \hline
    Leituras com data no futuro eram aceitas; data de recebimento informada pelo chamador & Dado com data falsa entrava no histórico \\ \hline
    Reenvio de lote já gravado gerava falha de sensor & O log registrava falhas que não aconteceram \\ \hline
    Todos os acessos com o superusuário \texttt{postgres} & Sem confidencialidade nem separação de funções \\ \hline
    Sem registro de quem alterou regras ou reconheceu eventos & Um limite de alarme podia ser elevado sem rastro \\ \hline
    Sem backup e recuperação & Uma falha de disco apagaria o histórico \\ \hline
  \end{tabularx}
  \fonte{Elaborado pelos autores (2026).}
\end{quadro}

\section{Arquitetura}

A arquitetura segue as camadas sensor, borda e nuvem de \textcite{grzesik2020}, com um servidor de aplicação entre o banco e os usuários (\autoref{fig:arquitetura}). Cada estação mede a cada minuto e envia um lote em JSON a cada dez minutos, identificado por uma chave própria. O servidor recebe os lotes, atende o painel, envia atualizações em tempo real e executa tarefas agendadas. O banco concentra as regras: valida cada leitura, gera eventos e alarmes, controla o acesso e registra a auditoria.

\begin{figure}[htb]
  \caption{Arquitetura do sistema}\label{fig:arquitetura}
  \centering
  \resizebox{0.92\textwidth}{!}{% Arquitetura: sensor -> borda -> servidor -> banco, com papéis de acesso
\begin{tikzpicture}[
  font=\footnotesize,
  caixa/.style={draw=azul, line width=0.7pt, rounded corners=3pt, fill=azulclaro, align=center, inner sep=4pt},
  sub/.style={draw=azul!60, rounded corners=2pt, fill=white, align=center, inner sep=2.5pt, font=\scriptsize, minimum width=2.45cm},
  disco/.style={cylinder, shape border rotate=90, aspect=0.22, draw=cinza, fill=cinzaclaro, align=center, font=\scriptsize, minimum width=1.9cm, minimum height=1.05cm},
  seta/.style={-{Stealth[length=2.2mm]}, line width=0.7pt, color=tinta},
  rot/.style={font=\scriptsize\itshape, color=tinta, align=center, fill=white, inner sep=1pt},
]
% --- campo
\node[caixa, minimum width=2.6cm, minimum height=1.3cm] (est1) at (0,1.1) {\textbf{Estação Centro}\\sensores + datalogger};
\node[caixa, minimum width=2.6cm, minimum height=1.3cm] (est2) at (0,-1.1) {\textbf{Estação Zona Norte}\\sensores + datalogger};
\node[font=\scriptsize\bfseries, color=cinza] at (0,2.35) {CAMPO (BORDA)};

% --- servidor
\node[draw=azul, line width=0.8pt, rounded corners=4pt, minimum width=3.4cm, minimum height=5.4cm, fill=azulclaro!40] (srv) at (4.6,0) {};
\node[font=\small\bfseries, anchor=north] at (srv.north) {Servidor Node.js};
\node[sub] (api) at (4.6,1.75) {API HTTP (JSON)\\\texttt{/api/ingest}};
\node[sub] (sse) at (4.6,0.75) {Tempo real (SSE)};
\node[sub] (age) at (4.6,-0.2) {Agendador};
\node[sub] (bkp) at (4.6,-1.15) {Backup (\texttt{pg\_dump})};
\node[sub] (cont) at (4.6,-2.1) {Fila de contingência};

% --- banco
\node[draw=verde, line width=0.8pt, rounded corners=4pt, minimum width=3.6cm, minimum height=5.4cm, fill=verdeclaro!50] (bd) at (9.3,0) {};
\node[font=\small\bfseries, anchor=north, align=center] at (bd.north) {PostgreSQL};
\node[sub, draw=verde!70] at (9.3,1.75) {Papéis + RLS};
\node[sub, draw=verde!70] at (9.3,0.8) {Validação e triggers};
\node[sub, draw=verde!70] at (9.3,-0.15) {Partições mensais};
\node[sub, draw=verde!70] at (9.3,-1.1) {Auditoria encadeada};
\node[sub, draw=verde!70] at (9.3,-2.05) {WAL (durabilidade)};

% --- usuários e discos
\node[caixa, minimum width=2.6cm, minimum height=1.1cm] (nav) at (4.6,-4.1) {\textbf{Navegador}\\painel HTML/CSS/JS};
\node[disco] (dbak) at (0.6,-3.0) {cópias\\.dump};
\node[disco] (dfila) at (0.6,-4.45) {fila em\\disco};

% --- ligações
\coordinate (ent) at (2.35,1.75);
\draw[line width=0.7pt, color=tinta] (est1.east) -- (ent);
\draw[line width=0.7pt, color=tinta] (est2.east) -- (ent);
\draw[seta] (ent) -- (api.west);
\node[rot, align=center] at (2.05,0.0) {lote JSON\\a cada 10 min\\+ chave da\\estação};
\draw[seta, <->] (srv.east |- 0,1.2) -- node[rot, above]{token da\\sessão} (bd.west |- 0,1.2);
\draw[seta] (bd.west |- 0,-0.55) -- node[rot, below]{NOTIFY} (srv.east |- 0,-0.55);
\draw[seta, <->] (nav.north) -- node[rot, right]{HTTPS\,/\,SSE} (srv.south);
\draw[seta, dashed] (bkp.west) -- ++(-0.95,0) |- (dbak.east);
\draw[seta, dashed] (cont.west) -- ++(-0.75,0) |- (dfila.east);
\end{tikzpicture}
}
  \fonte{Elaborado pelos autores (2026), adaptado de \textcite{grzesik2020}.}
\end{figure}

A decisão central é que o banco \textbf{não confia na aplicação}: o servidor se conecta com papéis de privilégio mínimo e repassa apenas o \textit{token} da sessão do usuário; é o banco que o valida e descobre quem é o usuário, seu perfil e as estações que pode ver. Mesmo com o servidor web comprometido, um invasor não se passa por outro usuário sem um \textit{token} válido.

\section{Metodologia de simulação e testes}

Um simulador gera leituras realistas (ciclo dia e noite, sazonalidade, tardes chuvosas e um dia de temporal), com cerca de 2\,\% de lotes atrasados e 0,5\,\% de leituras com o código de erro $-9999$. A carga cobre seis meses; depois, o servidor continua transmitindo em tempo real. Os testes foram de segurança (tentar burlar cada regra, inclusive como superusuário), de continuidade (derrubar o banco, fazer e restaurar backups) e de desempenho, com a metodologia de \textcite{grzesik2020}: escrita de 1.000 pontos, um por envio e em lotes de dez, e consultas de agregação repetidas 500 vezes em quatro cenários, antes e depois das melhorias, no mesmo ambiente.
